\documentclass[12pt,a4paper]{article}

% --- Настройка языка и шрифтов для LuaLaTeX ---
\usepackage{fontspec}
\defaultfontfeatures{Ligatures=TeX,Scale=MatchLowercase}

\setmainfont{Alice-Regular}[
    Path=font/,
    Extension=.ttf,
    BoldFont = Alice-Regular,
    ItalicFont = Alice-Regular,
    BoldItalicFont = Alice-Regular,
    BoldFeatures = {FakeBold=1.6},
    ItalicFeatures = {FakeSlant=0.2},
    BoldItalicFeatures = {FakeBold=1.6,FakeSlant=0.2}
]

\newfontfamily\cyrillicfont{Alice-Regular}[
    Path=font/,
    Extension=.ttf,
    BoldFont = Alice-Regular,
    ItalicFont = Alice-Regular,
    BoldItalicFont = Alice-Regular,
    BoldFeatures = {FakeBold=1.6},
    ItalicFeatures = {FakeSlant=0.2},
    BoldItalicFeatures = {FakeBold=1.6,FakeSlant=0.2}
]

\usepackage{polyglossia}
\setmainlanguage{russian}
\setsansfont{TeX Gyre Heros}
\setmonofont{TeX Gyre Cursor}
\AtBeginDocument{\renewcommand{\textnumero}{{\sffamily №}}}

\usepackage{graphicx}
\usepackage{amsmath, amsfonts, amssymb, amsthm}
\usepackage{mathtools}
\usepackage[protrusion=false]{microtype}
\usepackage{unicode-math}
\usepackage{array}
\usepackage{float}
\usepackage{caption}
\usepackage{enumitem}

\usepackage[paper=letterpaper,
            top=20mm,
            bottom=20mm,
            left=10mm,
            right=10mm,
            footskip=0.4in 
]{geometry}

% --- Колонтитулы ---
\usepackage{fancyhdr}
\pagestyle{fancy}
\fancyhf{}
\setlength{\headheight}{15pt}

\renewcommand{\headrulewidth}{0.4pt}
\renewcommand{\footrulewidth}{0.4pt}

\newcommand{\DocNumber}{2}
\newcommand{\DocTitle}{Перспективы глобальной трансформации 
техногенной цивилизации при тысячекратном увеличении плотности 
энергии компактных источников тока}
\newcommand{\Discipline}{ЭВМ и периферийные устройства}

\title{
	Министерство науки и высшего образования 
    
    Российской Федерации
    
    Национальный исследовательский университет «МИЭТ»
    
    Институт интегральной электроники имени академика 
    
    К. А. Валиева

	\vspace{2em}
}
\author{
}
\date{}

\begin{document}

\maketitle
\thispagestyle{fancy}

\begin{center}
	\Large Эссе \textnumero~$\DocNumber$\\
	по дисциплине «\Discipline»\\
	«\DocTitle»
\end{center}

\vspace{15em}

\begin{flushright}
	Выполнил студент группы ЭН-$35$\\
	Андреев И. А.
\end{flushright}

\cfoot{Москва 2026}

\pagebreak

\fancyhead[L]{\DocTitle}
\fancyhead[R]{\Discipline}
\fancyfoot[C]{\thepage}

\begin{center}
    \Large\textbf{Эссе \textnumero~$\DocNumber$}\\
    \large\textbf{\DocTitle}
\end{center}

\section{Введение}

Увеличение плотности энергии компактных химических источников тока в 
$1000$ раз представляет собой фундаментальный цивилизационный разрыв.  
Современные литий-ионные аккумуляторы коммерческого сегмента обладают 
гравиметрической плотностью порядка 
$250\text{--}300\text{ Вт}\cdot\text{ч/кг}$ 
($0{,}9\text{--}1{,}08\text{ МДж/кг}$) при объемной плотности около 
$600\text{--}700\text{ Вт}\cdot\text{ч/л}$. Стандартный аккумулятор 
современного смартфона массой около $60\text{ г}$ накапливает от $15$ 
до $18\text{ Вт}\cdot\text{ч}$ 
($0{,}015\text{--}0{,}018\text{ кВт}\cdot\text{ч}$) электроэнергии. 
Увеличение этой емкости в $1000$ раз при сохранении массы и 
геометрических габаритов означает скачок удельной энергоемкости до 
$250\text{--}300\text{ кВт}\cdot\text{ч/кг}$ 
($900\text{--}1080\text{ МДж/кг}$). 

В объеме карманного устройства 
оказывается сконцентрировано более $17\text{ кВт}\cdot\text{ч}$ 
электроэнергии, что превосходит суточное потребление домохозяйства 
в развитых странах и более чем в $20$ раз превышает удельную 
теплоту сгорания авиационного керосина или бензина 
($\approx 43\text{--}46\text{ МДж/кг}$, 
или $12\text{--}13\text{ кВт}\cdot\text{ч/кг}$). 

Масштабирование 
такой технологии на сопутствующие отрасли ликвидирует физические 
ограничения на перенос и концентрацию энергии, перестраивая 
транспорт, энергетические сети.

\section{Полная декарбонизация транспортного сектора и преодоление 
барьера авиации}

По данным Международного энергетического агентства (IEA) и 
Шестого оценочного доклада Межправительственной группы экспертов 
по изменению климата (IPCC AR6 WGIII), на транспортный сектор 
приходится около $15\%$ всех глобальных антропогенных парниковых 
выбросов и порядка $23\text{--}24\%$ мировых выбросов $\text{CO}_2$, 
связанных с энергетикой. До сих пор тяжелый транспорт, морское 
судоходство и дальнемагистральная авиация оставались наименее 
поддающимися декарбонизации отраслями 
именно в силу недостаточной плотности энергии аккумуляторных батарей.

\subsection{Преодоление критического порога авиации}

В исследованиях группы Венкатасубраманиана Вишванатана и Александра 
Биллса «Performance Metrics Required of Next-Generation Batteries to 
Electrify Commercial Aircraft» (ACS Energy Letters, $2020$) 
сформулированы физико-математические критерии электрификации 
коммерческой авиации. Расчеты показывают, что коммерческий 
пассажирский авиалайнер класса Boeing $737$ или Airbus A$320$ требует 
для обеспечения экономически обоснованной дальности полета плотности 
энергии накопителя не менее $800\text{ Вт}\cdot\text{ч/кг}$ на уровне 
упакованной системы. При современных показателях в 
$250\text{--}300\text{ Вт}\cdot\text{ч/кг}$ масса батарейного блока 
превышает допустимую взлетную массу воздушного судна, вытесняя 
полезную нагрузку и пассажиров.

В соответствии с уравнением дальности полета Бреге, керосиновые 
двигатели обладают критическим преимуществом: масса самолета 
непрерывно уменьшается по мере выгорания топлива в крейсерском 
полете, тогда как масса электрохимической батареи остается постоянной 
от взлета до посадки.

При увеличении емкости в $1000$ раз удельная 
плотность блока достигает $250\text{ кВт}\cdot\text{ч/кг}$, что 
нивелирует балластный эффект. Аккумулятор массой всего 
$1{,}5\text{--}2\text{ тонны}$ способен заменить 
$70\text{--}90\text{ тонн}$ керосина, обеспечивая 
трансконтинентальные перелеты тяжелых лайнеров с нулевыми 
прямыми выбросами, устраняя не только выбросы углекислого газа, 
но и эмиссию оксидов азота и сажевых следов в верхних слоях атмосферы.

\subsection{Переход к повсеместному применению аэромобилей}

В работе Шашанка Срипада и Венкатасубраманиана Вишванатана 
«Performance Metrics Required of Next-Generation Batteries 
to Electrify Vertical Takeoff and Landing (VTOL) Aircraft» 
(ACS Energy Letters, $2018$) доказано, что вертикальный взлет и 
зависание требуют экстремальных пиковых мощностей. В существующих 
реалиях аппараты eVTOL ограничены дальностью полета в несколько 
десятков километров, при этом резерв массы на полезную нагрузку 
остается минимальным. 

При тысячекратном скачке емкости компактный 
блок массой $10\text{ кг}$ аккумулирует около 
$2{,}5\text{ МВт}\cdot\text{ч}$ энергии. Это позволяет летательным 
аппаратам находиться в воздухе непрерывно на протяжении нескольких 
суток без дозаправки, превращая персональный и общественный 
авиатранспорт в доминирующий вид городской мобильности. Происходит 
масштабная трансформация городского пространства: градостроительство 
освобождается от необходимости расширения плоскостных автомагистралей, 
развязок и парковочных площадей в пользу воздушных эшелонов и 
посадочных терминалов на крышах зданий.

\section{Деструкция централизованных энергосетей и автономная электрификация}

Инфраструктура современной энергетики базируется на парадигме 
централизованной генерации, где гигантские электростанции соединены с 
конечными потребителями разветвленной сетью магистральных ЛЭП, 
распределительных кабелей и трансформаторных подстанций. Эта топология 
сопряжена с высокими капитальными затратами, потерями при передаче 
энергии и уязвимостью перед масштабными блэкаутами.

Аналитический центр 
Rocky Mountain Institute (RMI) в исследовании Питера Бронски, Джона 
Крейтса и Леи Гуччионе «The Economics of Grid Defection» (2014) 
математически обосновал феномен «смертельной спирали коммунальных 
энергосетей». Механизм дефекции заключается в 
том, что по мере удешевления распределенной генерации и систем 
накопления энергии наиболее состоятельные домохозяйства и коммерческие 
предприятия отключаются от центральной сети. Это вынуждает регуляторов 
и коммунальные компании распределять фиксированные эксплуатационные 
затраты на содержание сотен тысяч километров электросетей среди 
уменьшающегося числа оставшихся абонентов, что ведет к неизбежному 
росту тарифов и стимулирует дальнейший лавинообразный уход 
потребителей.

Тысячекратное увеличение емкости доводит эту тенденцию 
до абсолютного предела. Накопитель размером с автомобильный 
аккумулятор и массой $3\text{--}5\text{ кг}$ аккумулирует от $750$ 
до $1400\text{ кВт}\cdot\text{ч}$ электроэнергии. Данного объема 
достаточно для полного энергообеспечения частного дома, включая 
климатические системы, в течение полугода автономного 
функционирования. Необходимость физического подключения жилья к 
централизованной энергосистеме исчезает: локальные солнечные панели 
или микротурбины заряжают компактный блок без привязки к балансу 
спроса и предложения в сети.

\subsection{Скачок в развитии государств Юга}

Как подчеркивается в отчете IRENA 
«A New World: The Geopolitics of the Energy Transformation» (2019), 
развивающиеся регионы Глобального Юга получают возможность совершить 
инфраструктурный скачок. Подобно тому как развивающиеся 
страны Африки и Азии миновали этап прокладки медных телефонных сетей, 
сразу внедрив сотовую связь, сотни миллионов жителей, лишенных доступа 
к электросетям, получают базовую энергоавтономию без необходимости 
многомиллиардных инвестиций в магистральные распределительные сети.

\subsection{Демонтаж рентной модели петрогосударств}

В исследовании IRENA (2019) показано, что переход к технологиям 
распределенной энергии лишает углеводородные автократии 
геополитической ренты. Когда личный, коммерческий и авиационный 
транспорт, а также отопительные системы полностью переходят на 
автономные накопители, мировой спрос на сырую нефть сокращается до 
объемов, востребованных исключительно химической промышленностью. Это 
влечет за собой бюджетный коллапс государств, экономика которых 
базируется на экспорте углеводородов, и ликвидирует стратегическое 
значение традиционных морских узлов транзита (Ормузский пролив, 
Малаккский пролив, Суэцкий канал).

\section{Физический дуализм: катастрофические риски концентрации 
энергии и режим безопасности}

Концентрация огромного потенциала в миниатюрном объеме порождает 
фундаментальную проблему физической безопасности. Закон сохранения 
энергии обуславливает неизбежный дуализм: любая среда высокой 
энергоемкости в случае катастрофической разгерметизации, внутреннего 
замыкания или механического разрушения ведет себя как мощное 
взрывчатое вещество.

\subsection{Термодинамический эквивалент детонации}

Удельная теплота взрыва стандартного тринитротолуола (ТНТ) 
составляет $4{,}184\text{ МДж/кг}$ ($4184\text{ Дж/г}$).

Расчет потенциальной энергии ячейки емкостью 
$17{,}325\text{ кВт}\cdot\text{ч}$ дает:
$$
E = 17{,}325\text{ кВт}\cdot\text{ч} \times 3{,}6 \times 10^6\text{ Дж/кВт}\cdot\text{ч} = 62{,}37 \times 10^6\text{ Дж} = 62{,}37\text{ МДж}
$$

Прямой тротиловый эквивалент мгновенно высвобождаемой энергии 
составляет:
$$
M_{\text{TNT}} = \frac{62{,}37\text{ МДж}}{4{,}184\text{ МДж/кг}} \approx 14{,}91\text{ кг ТНТ}
$$
Экспериментальные исследования теплового разгона современных 
аккумуляторов типоразмера 18650 емкостью около 
$10\text{ Вт}\cdot\text{ч}$ показывают, что в момент термической 
деструкции ячейка способна выделиться с энергией взрывной волны, 
эквивалентной $5\text{ граммам}$ тротила. При увеличении концентрации 
энергии в 1000 раз карманный телефон массой $180\text{--}200\text{ г}$ 
с батареей в $60\text{ г}$ приобретает разрушительный потенциал 
тяжелого артиллерийского фугасного снаряда калибра 
$152\text{--}155\text{ мм}$ или авиационной бомбы.

\subsection{Нормативный кризис гражданского оборота и правил авиаперевозок}

Международная организация гражданской авиации (ICAO Doc 9284) и 
Федеральное управление гражданской авиации США (FAA) вводят жесткие 
регуляторные рамки для портативных накопителей:Аккумуляторы емкостью 
до $100\text{ Вт}\cdot\text{ч}$ допускаются к свободной перевозке 
пассажирами в ручной клади.Аккумуляторы емкостью от $101$ до 
$160\text{ Вт}\cdot\text{ч}$ требуют обязательного согласования с 
эксплуатантом воздушного судна.Любые накопители емкостью свыше 
$160\text{ Вт}\cdot\text{ч}$ категорически запрещены к перевозке на 
пассажирских бортах.Батарея смартфона емкостью 
$17\,325\text{ Вт}\cdot\text{ч}$ превышает порог ICAO более чем в 
108 раз. В случае короткого замыкания в салоне произойдет разрушение 
силового контура планера и гарантированная гибель воздушного судна.

В результате свободное обращение потребительской электроники в 
гражданском обществе станет юридически невозможным. Возникнет 
необходимость проектирования тяжелых композитных взрывозащитных 
оболочек, введения государственной системы лицензирования и номерного 
учета каждого аккумулятора аналогично огнестрельному оружию, 
радиоактивным материалам и взрывчатке. Городские пространства, 
метрополитены, правительственные учреждения и общественный транспорт 
будут вынуждены перейти на тотальный досмотр с помощью рентгеновской 
и спектрометрической аппаратуры для предотвращения террористических 
актов с применением бытовых устройств.

\section{Заключение}

Историко-технологический анализ смены цивилизационных укладов, 
выполненный Вацлавом Смилом в трудах «Energy and Civilization: 
A History» ($2017$) и «Energy Transitions: History, Requirements, 
Prospects» ($2010$), демонстрирует, что динамика человеческого 
развития определялась переходом к энергоносителям со все более 
высокой концентрацией энергии на единицу массы: от сухой биомассы 
($15\text{ МДж/кг}$) к каменному углю ($25\text{ МДж/кг}$) и 
сырой нефти ($42\text{--}44\text{ МДж/кг}$). Доминирование жидких 
углеводородов на протяжении XX и XXI веков сохранялось именно 
благодаря непревзойденной плотности энергии, недостижимой для 
химических аккумуляторов, удельная емкость которых со времен 
свинцово-кислотных элементов $1900$ года 
($25\text{ Вт}\cdot\text{ч/кг}$) выросла лишь на порядок.

Появление 
накопителя с показателем порядка $1000\text{ МДж/кг}$ 
($280\text{ кВт}\cdot\text{ч/кг}$) знаменует преодоление теплотворной 
способности любого химического топлива более чем в двадцать раз, что 
влечет за собой глобальное переворот мироустройства. Транспорт, 
энергетика, социальный уклад и геополитическая конфигурация государств
предмет радикальной трансформации.

\end{document}
